\section{Enterprise Scaling and Sizing}

This section addresses the question of what deploying the system at
organisational scale actually requires. Figures marked \emph{measured} derive
from the reference implementation; those marked \emph{projected} are
extrapolations and are stated as such.

\subsection{The sizing model}

Three quantities determine every hardware decision.

\begin{align*}
\text{chunks} &= \text{pages} \times 3.3 && \text{\small(measured)}\\
\text{index bytes} &= \text{chunks} \times 12\,\text{KB}
   && \text{\small(5.4\,KB vectors + 6.2\,KB graph)}\\
\text{ingest GPU-hours} &= \frac{\text{chunks} \times t_{\text{chunk}}}{3600}
   && \text{\small($t_{\text{chunk}} \approx 15$\,s optimised, projected)}
\end{align*}

\noindent Per-chunk storage decomposes as: a 1024-dimensional
\code{float32} embedding (4\,096\,B), the chunk payload and metadata
($\approx$1\,300\,B), approximately 8.5 graph nodes and 6.6 edges
($\approx$6.2\,KB). Source documents are additional and typically dominate raw
storage while being irrelevant to query performance.

\subsection{Worked storage examples}

\begin{center}
\renewcommand{\arraystretch}{1.25}
\begin{tabular}{@{}l r r r r@{}}
\toprule
\textbf{Corpus} & \textbf{Pages} & \textbf{Chunks} & \textbf{Index} & \textbf{Ingest (1 GPU)} \\
\midrule
Departmental archive & 50\,000 & 165\,K & 2.0\,GB & $\approx$29 days \\
Plant document set & 500\,000 & 1.65\,M & 20\,GB & $\approx$10 months \\
Refinery, all history & 5\,000\,000 & 16.5\,M & 198\,GB & impractical \\
Enterprise, multi-site & 20\,000\,000 & 66\,M & 792\,GB & impractical \\
\bottomrule
\end{tabular}
\end{center}

\begin{keybox}{Storage is cheap; ingestion is the real constraint}
Index storage is modest even at twenty million pages --- under a terabyte, well
within a single NVMe array. The binding constraint is \emph{GPU-hours for
extraction}. Bulk ingestion at plant scale is a parallel batch problem
requiring a dedicated ingestion cluster, not a bigger disk.
\end{keybox}

\subsection{Mitigating the ingestion cost}

\begin{description}[leftmargin=0pt,style=nextline]
\item[Parallelism.] Extraction is embarrassingly parallel across chunks.
$N$ GPUs deliver close to $N\times$ throughput; 16 GPUs reduce the 500\,000-page
case from ten months to roughly nineteen days.
\item[Continuous batching.] Replacing Ollama with vLLM for the extraction
worker enables batched decoding, empirically worth a further 4--8$\times$ on
short structured outputs.
\item[Tiered extraction.] Not every document justifies graph extraction.
Indexing everything for vector search while restricting graph extraction to
governing documents --- procedures, drawings, incident reports --- typically
reduces the graph workload by an order of magnitude.
\item[Incremental ingestion.] After the initial load, only new and revised
documents are processed. Steady-state volume is a small fraction of the bulk
load.
\end{description}

\subsection{Deployment tiers}

\begin{center}
\small
\renewcommand{\arraystretch}{1.3}
\begin{tabularx}{\linewidth}{@{}l l l l r X@{}}
\toprule
\textbf{Tier} & \textbf{GPU} & \textbf{RAM} & \textbf{Storage} & \textbf{Users} & \textbf{Corpus} \\
\midrule
Pilot & 1 $\times$ 16--24\,GB & 64\,GB & 2\,TB NVMe & 5--15 & $\le$50\,K pages \\
Departmental & 1 $\times$ 48\,GB & 128\,GB & 8\,TB NVMe & 25--75 & $\le$500\,K pages \\
Plant-wide & 2--4 $\times$ 80\,GB & 256--512\,GB & 30\,TB NVMe & 100--300 & $\le$5\,M pages \\
Enterprise & 8--16 $\times$ 80\,GB & 1\,TB+ & 100\,TB+ & 1\,000+ & 20\,M+ pages \\
\bottomrule
\end{tabularx}
\end{center}

\begin{center}
\small
\renewcommand{\arraystretch}{1.3}
\begin{tabularx}{\linewidth}{@{}l X@{}}
\toprule
\textbf{Tier} & \textbf{Characteristics} \\
\midrule
\textbf{Pilot} & Exactly the reference implementation. One model resident at a
time with eviction. Single workstation under a desk; no data-centre
dependency. Proves value before procurement. \\
\addlinespace
\textbf{Departmental} & A single 48\,GB card holds the reasoning and vision
models concurrently, removing swap latency. Supports the full \code{venue}
profile without eviction. \\
\addlinespace
\textbf{Plant-wide} & 80\,GB cards enable 120B-class open-weight models,
matching the problem statement's \code{datacenter} profile. One GPU is
dedicated to batch ingestion; the remainder serve interactive load. Graph and
vector stores migrate from embedded to clustered. \\
\addlinespace
\textbf{Enterprise} & Separate ingestion and serving clusters, high
availability, per-site corpora with federated retrieval, role-based access
control over graph partitions. \\
\bottomrule
\end{tabularx}
\end{center}

\subsection{Concurrency}

Interactive capacity is governed by key-value cache memory, not by model
weights:

\[
\text{concurrent sessions} \approx
\frac{V_{\text{GPU}} - V_{\text{weights}}}{V_{\text{KV per session}}}
\]

\noindent For a 14B model at 4-bit ($\approx$9\,GB) on an 80\,GB card with
roughly 1\,GB of KV cache per 8\,K-token session, the arithmetic gives about 70
sessions; 40--50 is a realistic planning figure once fragmentation and headroom
are accounted for. Because agent runs are bursty rather than sustained, a
served-user multiple of five to ten times the concurrent-session count is
typical.

\subsection{Architectural changes required above the pilot tier}

The reference implementation makes deliberate choices that are correct for a
single workstation and must change at scale.

\begin{center}
\renewcommand{\arraystretch}{1.25}
\begin{tabularx}{\linewidth}{@{}l l X@{}}
\toprule
\textbf{Component} & \textbf{Pilot} & \textbf{Scale-out replacement} \\
\midrule
Graph store & K\`uzu (embedded) & Neo4j cluster or K\`uzu with a read-replica fan-out \\
Vector store & Qdrant local & Qdrant or Milvus cluster with sharding \\
Inference & Ollama & vLLM with continuous batching and tensor parallelism \\
Cache & SQLite & PostgreSQL or Redis \\
Sandbox & \code{unshare} & gVisor or Kata Containers under an orchestrator \\
Identity & none & OIDC with role-based graph partition access \\
Ingestion & in-process & a queue with distributed workers \\
\bottomrule
\end{tabularx}
\end{center}

\begin{keybox}{A caution on the single-writer constraint}
K\`uzu permits one writing process. In the reference implementation this means
indexing and serving cannot run concurrently --- an acceptable limitation for a
pilot and an unacceptable one for production. It is the first component that
must be replaced when moving beyond a single workstation.
\end{keybox}
